\begin{frame}{Onde isso aparece}
  \begin{columns}[T,onlytextwidth]
    \begin{column}{0.49\textwidth}
      \begin{destaque}{Moléculas diatômicas}
        Vibrações em torno do equilíbrio. As transições
        $\Delta n = \pm 1$ explicam as bandas no infravermelho.
      \end{destaque}
      \begin{destaque}[fisicaamber]{Íons aprisionados}
        O movimento de um íon em um potencial efetivo forma um
        oscilador; seus estados são usados em computação quântica.
      \end{destaque}
    \end{column}
    \begin{column}{0.49\textwidth}
      \begin{destaque}[fisicaviolet]{Circuitos supercondutores}
        Junções Josephson formam um oscilador anarmônico cujos dois
        primeiros níveis definem um qubit.
      \end{destaque}
      \begin{destaque}[fisicagreen]{Nanomecânica}
        Osciladores mecânicos levados ao estado fundamental testam
        os limites da mecânica quântica.
      \end{destaque}
    \end{column}
  \end{columns}
\end{frame}
